% Capítulo 6: Conclusiones y trabajo futuro.

Este capítulo recapitula lo realizado y lo observado en relación con
los objetivos planteados al inicio. La primera sección sintetiza las
conclusiones del trabajo, distinguiendo lo que sostienen los resultados
de lo que queda como interpretación prudente. La segunda recoge las
líneas de trabajo futuro derivadas de las limitaciones detectadas.

\section{Conclusiones}

El trabajo ha construido una referencia clásica rigurosa de selección de
características y la ha utilizado para evaluar un método QFS basado en
átomos neutros. La conclusión principal no es que QFS gane o pierda en
un marcador agregado, sino que la referencia clásica permite leer el
proceso: primero caracteriza el régimen del dato, después sitúa cada
selector en el compromiso relevancia--redundancia y, finalmente,
permite separar si un deterioro procede del criterio QUBO o de la
optimización analógica.

El primer resultado es metodológico. La cadena experimental de diez
fases deja trazado cada eslabón: análisis del dato crudo,
preprocesado conservador, comprobación posterior al preprocesado, particiones sin
solapes ni \textit{leakage} detectable, selección clásica con doce
métodos, modelado con contrastes, comparación final, ejecución QFS,
integración cuántica y síntesis visual. Esta estructura evita que la
fase cuántica se evalúe sobre datos sin contexto. La corrección FDR, los
tamaños de efecto, los VIF, la PCA, el \textit{drift}, la validación
adversarial y las permutaciones no son controles accesorios: son los que
permiten interpretar después cada resultado de QFS como consecuencia de
un régimen medido.

Sobre la pregunta central (si el método cuántico mejora o empeora a la
selección clásica), la conclusión de esta evaluación es directa y
prudente: QFS \emph{no supera} a la selección de características clásica.
En los regímenes con señal clara (\textit{Breast Cancer} y las dos
formulaciones de \textit{Olive Oil}) iguala tanto a la referencia clásica
como al uso de todas las variables, sin mejorarlos de forma
significativa; en los regímenes difíciles (\textit{Madelon} y
\textit{Customer Churn}) se deteriora. La selección
clásica, en cambio, iguala o mejora al uso de todas las variables en
todas las formulaciones y aporta una ganancia sustancial cuando abundan
los distractores. En los conjuntos y el protocolo evaluados, por tanto,
QFS es en el mejor de los casos competitivo con la selección clásica,
pero no la aventaja. Este resultado negativo parcial no resta valor al
estudio, sino que delimita con precisión el alcance del método: QFS no
aparece como un sustituto general de la selección clásica, sino como una
formulación cuyo interés depende del régimen del dato y de la separación
entre criterio y optimización.

El resultado más informativo, sin embargo, no es ese veredicto agregado,
sino su descomposición. Comparar el método cuántico con el óptimo exacto
de su propia función objetivo revela que, cuando QFS se deteriora, el control exacto distingue dos
mecanismos según el régimen, en lugar de una regla general. En
\textit{Madelon}, el deterioro procede del criterio (la relevancia y la
redundancia estimadas mediante información mutua entre pares de variables),
que no captura la señal cuando esta depende de interacciones de orden
superior: el óptimo exacto ya queda lejos del baseline. En \textit{Customer
Churn}, en cambio, el criterio casi alcanza al baseline y la pérdida se
concentra en la lectura analógica de densidades (el \textit{readout}), no
en el criterio. El cuello de botella, por tanto, no es único: en los
regímenes de señal escasa es de formulación del criterio, mientras que en
los de baja redundancia es de lectura del optimizador analógico. La vía de
mejora se diagnostica por separado en cada caso. Los detalles por conjunto
que sostienen esta lectura (qué variables retiene cada subconjunto, dónde se
concentra el deterioro y por qué) se desarrollan en el capítulo de
resultados y discusión.

El resultado final responde a los objetivos del TFG\@. Se ha revisado y
ejecutado una referencia clásica amplia; se han seleccionado datasets
que cubren regímenes distintos; se ha comparado el rendimiento de los
selectores bajo un protocolo común; se ha implementado y evaluado QFS
conservando el núcleo del método de referencia y declarando sus
adaptaciones operativas; y se ha comparado frente a los clásicos con
una lectura atribuible. La aportación específica es metodológica: el
bloque clásico se convierte en un instrumento para predecir por régimen y
atribuir por mecanismo, más allá de su papel de referencia de
rendimiento.

Más allá del diagnóstico, que QFS pueda igualar a la selección clásica en
varios regímenes tiene una lectura computacional relevante de cara al
hardware real. La formulación sobre átomos neutros asigna un átomo (un
cúbit) a cada variable candidata, de modo que un dispositivo de $n$
cúbits puede seleccionar características sobre un espacio de hasta $n$
variables resolviendo el compromiso relevancia--redundancia de forma
global y en un tiempo de ejecución muy corto, frente al coste
combinatorio de la búsqueda exacta clásica. Las plataformas industriales
de átomos neutros más habituales operan ya con cientos de cúbits
accesibles en la nube (el procesador Aquila de QuEra, de 256 cúbits, está
disponible sobre Amazon Braket~\cite{wurtz2023}) y han superado el millar
de átomos atrapados en un solo disparo en laboratorio~\cite{pasqal2024},
con demostraciones de cúbits lógicos corregidos sobre esta misma
tecnología~\cite{bluvstein2024}. El escalado previsto apunta a los miles
de cúbits en los próximos años, lo que permite considerar esta línea como
una proyección experimental plausible, aunque su viabilidad práctica
dependerá de las restricciones de ruido, conectividad efectiva,
calibración, escalado del embebido y coste de repetición de medidas. Que el criterio resulte competitivo es, por
tanto, la condición necesaria para que esa ventaja de eficiencia merezca
llevarse a hardware físico, y es lo que dota de sentido práctico a la
proyección que se plantea a continuación.

\section{Trabajo futuro}

La primera línea futura afecta al criterio. \textit{Madelon} muestra que
una formulación basada en información mutua por pares puede quedarse
corta cuando la señal depende de interacciones de orden superior. Una
extensión natural es estudiar criterios de relevancia conjunta,
términos de orden mayor o formulaciones HUBO (\textit{Higher-order
Unconstrained Binary Optimization}) que permitan capturar esas
interacciones antes de llevar el problema al optimizador cuántico. Esta
dirección está alineada con la evolución natural de la formulación QUBO,
pero debe evaluarse con el mismo tipo de control frente al óptimo exacto.

La segunda línea afecta al readout y a la dinámica analógica en regímenes
de baja redundancia. \textit{Customer Churn} no apunta a un fallo
geométrico del MDS, pero sí a un déficit leve del optimizador cuando el
criterio es casi correcto y el baseline está saturado. Un estudio
dirigido podría analizar variantes de lectura de densidades, protocolos
de \textit{driving}, escalado de distancias o presupuestos de shots para
comprobar si el hueco frente al oráculo se reduce.

La tercera línea es una ablación de preprocesado, acotada a
\textit{Customer Churn}. El \textit{one-hot} utilizado en este trabajo es
defendible para variables nominales y para los modelos clásicos, pero
crea grupos de dummies que compiten por el presupuesto de QFS\@. Una
comparación específica entre \textit{one-hot}, \textit{one-hot}
\texttt{drop\_first} y \textit{label encoding} permitiría separar el
efecto de la codificación del régimen intrínseco de baja redundancia, sin
reabrir todo el banco experimental.

La cuarta línea es la consistencia fuerte. El estudio actual mide
estabilidad de selección con tres semillas y robustez estadística por
bootstrap y permutación; QFS usa 100 inicializaciones MDS, pero no se ha
realizado un experimento multi-semilla completo de modelado y QFS\@. Si el
coste computacional lo permite, varias reejecuciones completas de fase
8--9 permitirían cuantificar la variabilidad del readout analógico y no
solo su error frente al óptimo exacto.

Por último, la ejecución sobre hardware real de átomos neutros queda
como proyección natural. La implementación evaluada opera en simulador,
como la referencia metodológica disponible, pero incorpora decisiones
operativas propias (preselección híbrida, codificación \textit{one-hot}
en \textit{Customer Churn} y barrido de $\beta$). El salto a hardware
físico introduciría ruido, calibración y
restricciones experimentales propias, y debería plantearse como una
extensión del estudio, no como una condición para validar las
conclusiones aquí obtenidas.
